\origpage{345}{N. Robertson, P.D. Seymour / Journal of Combinatorial Theory, Series B 90 (2004) 325--385}
\hspace*{12pt}$\bullet$ $ins(C^*) \subseteq \bar{\Lambda}$ for every circuit $C^*$ of $\Gamma^*$ with $U(C^*) \subseteq \bar{\Lambda}$ and $|E(C^*)| < 2(rep(\mathscr{T}) - \kappa - 5)$.

\textbf{Proof.} For $1 \le i \le rep(\mathscr{T})$, let $Z(i)$ be the union of $ins(W)$, taken over all closed walks $W$ of $\Gamma^*$ with length $< 2i$ and with $z^* \subseteq ins(W)$, where $z^*$ is the atom of $\Gamma^*$ corresponding to $z$. By Theorems 8.5, 8.10, 8.12 and 9.2 of [6], $Z(rep(\mathscr{T}))$ is simply-connected and $\ne \hat{\Sigma}$, and so by Theorems 4.2 and 11.9 of [2], there is a closed disc $\Delta \subseteq \hat{\Sigma}$ with $Z(rep(\mathscr{T})) \subseteq \Delta \setminus bd(\Delta)$. Since $\kappa \ge 4$, it follows that $z \subseteq Z(\kappa) \subseteq Z(rep(\mathscr{T})) \subseteq \Delta$.

\hspace*{12pt}(1) \textit{Let $r$ be a region of $sk(\Gamma)$ in $\hat{\Sigma}$ with $\bar{r} \cap Z(\kappa) \ne \emptyset$, and let $s \in A(\Gamma)$ with $r \subseteq s$; then $d(z, s) \le \kappa + 2$ and $s \subseteq \Delta \setminus bd(\Delta)$.}

\textit{Subproof.} Since $\bar{r} \cap Z(\kappa) \ne \emptyset$ it follows that $\bar{s} \cap Z(\kappa) \ne \emptyset$, and so there is a vertex $v$ of $\Gamma$ with $v \in \bar{s} \cap Z(\kappa)$. Hence $d(s, v) \le 3$ and $d(z, v) < \kappa$, and so $d(z, s) \le \kappa + 2$. Suppose that $s \nsubseteq Z(rep(\mathscr{T}))$. Since either $r = s$ or $s \in E(\Gamma)$ with $|\bar{s}| = 3$, $s$ is a subset of the closure of the union of the regions of $\Gamma^*$ that meet $s$; and so there is a region $e^*$ of $\Gamma^*$ with $s \cap e^* \ne \emptyset$ and $e^* \nsubseteq Z(rep(\mathscr{T}))$. Let $e$ be the corresponding edge of $\Gamma$; then $d(z, e) = rep(\mathscr{T})$ since $e^* \nsubseteq Z(rep(\mathscr{T}))$. But $d(s, e) \le 3$ since $s \cap e^* \ne \emptyset$, and so
\begin{equation*}
d(z, s) \ge rep(\mathscr{T}) - 3 \ge \kappa + 3,
\end{equation*}
a contradiction. Thus $s \subseteq Z(rep(\mathscr{T})) \subseteq \Delta \setminus bd(\Delta)$. This proves (1).

From (1) and Theorem 5.2 of [8] it follows that there is a circuit $C$ of $sk(\Gamma)$ with $U(C) \subseteq \Delta$, bounding an open disc $\Lambda \subseteq \Delta$ including $Z(\kappa)$, such that every edge of $C$ is incident with a region $r$ of $sk(\Gamma)$ in $\hat{\Sigma}$ with $\bar{r} \cap Z(\kappa) \ne \emptyset$. We claim that $C$ satisfies the theorem. Certainly $z \subseteq \Lambda$, since $z \subseteq Z(\kappa)$. To verify the first assertion of the theorem, we shall show that $\bar{\Lambda} \subseteq Z(\kappa + 6)$. Let $e \in E(C)$, and let $r$ be a region of $sk(\Gamma)$ in $\hat{\Sigma}$ with $e \subseteq \bar{r}$ and with $\bar{r} \cap Z(\kappa) \ne \emptyset$. Let $s \in A(\Gamma)$ with $r \subseteq s$. Then by (1), $d(z, s) \le \kappa + 2$. Let $f \in E(\Gamma)$ with $e \subseteq f$; then $e \subseteq \bar{s} \cap f$ and so $d(s, f) \le 3$. Hence $d(z, f) \le \kappa + 5$, and so $f \subseteq f^* \subseteq Z(\kappa + 6)$, where $f^* \in A(\Gamma^*)$ corresponds to $f$. Consequently $e \subseteq Z(\kappa + 6)$. Since this holds for every edge $e$ of $C$ it follows that $U(C) \subseteq Z(\kappa + 6)$. By Theorem 8.10 of [6], $Z(\kappa + 6)$ is simply-connected, and so there is a closed disc $D \subseteq Z(\kappa + 6) \subseteq \Delta$ bounded by $U(C)$. Thus $D$ and $\bar{\Lambda}$ are both closed discs in $\Delta$ bounded by $U(C)$, and hence $D = \bar{\Lambda}$; and consequently $\bar{\Lambda} \subseteq Z(\kappa + 6)$. Now let $x \in A(\Gamma)$ with $x \cap \bar{\Lambda} \ne \emptyset$. Then $x \cap Z(\kappa + 6) \ne \emptyset$, and so $d(z, x) \le \kappa + 5$. Hence the first assertion of the theorem holds.

To verify the second, let $x \in A(\Gamma)$ with $x \cap \Lambda = \emptyset$, and let $x^*$ be the corresponding atom of $\Gamma^*$. Since $x \cap x^* \ne \emptyset$ and $x \cap \Lambda = \emptyset$ it follows that $x^* \nsubseteq \Lambda$, and so $x^* \nsubseteq Z(\kappa)$. Hence $d(z, x) \ge \kappa$. This verifies the second assertion.

For the third, let $x \in A(\Gamma)$ with $x \nsubseteq \Lambda$. If $x \cap \Lambda = \emptyset$ then because the second assertion of the theorem holds, $d(z, x) \ge \kappa$ as required. We may assume then that $x \cap \Lambda \ne \emptyset$, and since $x \nsubseteq \Lambda$ it follows that $x \cap U(C) \ne \emptyset$. Since $x \in A(\Gamma)$ and $x \cap \Lambda \ne \emptyset$ and $x \cap U(C) \ne \emptyset$, it follows that $x$ is an edge of $\Gamma$ with $|\bar{x}| = 3$, and there is a vertex
